% Telos shared preamble.
%
% Every Telos publication is designed to be printed on a home laser printer,
% one sheet at a time, and carried into a boat or a garage. Two constraints
% follow from that and govern everything below:
%
%   1. Pure monochrome. No value is ever a hue. Emphasis is carried by rule
%      weight, fill, and type, never by color, so a grayscale or black-only
%      printer loses nothing.
%   2. Card density. A "one-pager" is a hard contract: the leaf must fit on a
%      single side of US Letter. The layout is therefore tight by default and
%      the environments below are sized to leave room for a plate.

\documentclass[10pt,letterpaper]{article}

\usepackage[margin=0.55in,includefoot,footskip=0.24in]{geometry}
\usepackage[T1]{fontenc}
\usepackage[utf8]{inputenc}
\usepackage{lmodern}
\usepackage{microtype}
\usepackage{array}
\usepackage{booktabs}
\usepackage{longtable}
\usepackage{tabularx}
\usepackage{enumitem}
\usepackage{needspace}
\usepackage{multicol}
\usepackage{ragged2e}
\usepackage{graphicx}
\usepackage{wrapfig}
\usepackage{xcolor}
\usepackage{fancyhdr}
\usepackage{titlesec}
\usepackage{hyperref}
\usepackage[most]{tcolorbox}
\usepackage{tikz}
\usetikzlibrary{arrows.meta,calc,positioning,shapes.geometric,decorations.pathmorphing,patterns,fadings}

% Semantic names are kept so a document reads intelligibly, but every value is
% black, white, or a neutral gray that survives a black-only print driver.
\definecolor{telosink}{gray}{0}
\definecolor{telospaper}{gray}{1}
\definecolor{telosrule}{gray}{0.35}
\definecolor{telosfill}{gray}{0.90}
\definecolor{telosdeep}{gray}{0.72}
\definecolor{telosmute}{gray}{0.45}

\hypersetup{hidelinks,pdfborder={0 0 0}}

% Keep an unchanged source byte-reproducible across rebuilds: no build-clock
% creation date and no random trailer ID.
\pdfinfoomitdate=1
\pdftrailerid{}

\setlength{\parindent}{0pt}
\setlength{\parskip}{0.42em}
\setlength{\tabcolsep}{4pt}
\setlength{\columnsep}{1.1em}
\renewcommand{\arraystretch}{1.18}
\setlist[itemize]{leftmargin=1.15em,itemsep=0.12em,topsep=0.2em,parsep=0pt}
\setlist[enumerate]{leftmargin=1.5em,itemsep=0.18em,topsep=0.2em,parsep=0pt}
\setlist[description]{leftmargin=0pt,itemsep=0.16em,topsep=0.2em,parsep=0pt,
  font=\normalfont\bfseries}

% ---------------------------------------------------------------------------
% Document identity. Each leaf declares these before \begin{document} so the
% running head, the colophon, and the PDF metadata all agree.
% ---------------------------------------------------------------------------
\newcommand{\TelosProject}{Telos}
\newcommand{\TelosDocTitle}{}
\newcommand{\TelosDocSubtitle}{}
\newcommand{\TelosRevision}{}

\newcommand{\telosmeta}[4]{%
  \renewcommand{\TelosProject}{#1}%
  \renewcommand{\TelosDocTitle}{#2}%
  \renewcommand{\TelosDocSubtitle}{#3}%
  \renewcommand{\TelosRevision}{#4}%
  \hypersetup{pdftitle={#2},pdfsubject={#3; version #4},
    pdfkeywords={Telos, version #4},pdfauthor={Telos}}%
}

\pagestyle{fancy}
\fancyhf{}
\renewcommand{\headrulewidth}{0pt}
\renewcommand{\footrulewidth}{0.4pt}
\fancyfoot[L]{\footnotesize\textsc{\TelosProject}}
\fancyfoot[C]{\footnotesize\TelosDocTitle}
\fancyfoot[R]{\footnotesize\thepage}

% The first page of a one-pager carries the masthead instead of a running foot,
% so its footer is the compact provenance line.
\fancypagestyle{telosfirst}{%
  \fancyhf{}%
  \renewcommand{\headrulewidth}{0pt}%
  \renewcommand{\footrulewidth}{0.4pt}%
  \fancyfoot[L]{\footnotesize\textsc{\TelosProject}}%
  \fancyfoot[C]{\footnotesize\TelosRevision}%
  \fancyfoot[R]{\footnotesize\thepage}%
}

% ---------------------------------------------------------------------------
% Masthead. A heavy rule, the title, a light subtitle, a thin rule. The
% optional argument is a right-aligned tag (a species' scientific name, a
% lesson number) that sits on the title's baseline.
% ---------------------------------------------------------------------------
\newcommand{\telosmasthead}[1][]{%
  \thispagestyle{telosfirst}%
  \begingroup
  \setlength{\parskip}{0pt}%
  \noindent\rule{\linewidth}{1.6pt}\par
  \vspace{0.30em}
  \noindent
  \begin{minipage}[b]{0.66\linewidth}
    \raggedright{\LARGE\bfseries\TelosDocTitle}
  \end{minipage}%
  \hfill
  \begin{minipage}[b]{0.32\linewidth}
    \raggedleft{\small\itshape #1}
  \end{minipage}\par
  \vspace{0.18em}
  \noindent{\small\TelosDocSubtitle}\par
  \vspace{0.28em}
  \noindent\rule{\linewidth}{0.4pt}\par
  \endgroup
  \vspace{0.35em}
}

% ---------------------------------------------------------------------------
% Headings. Compact, small-caps, rule-anchored — a card has no room for the
% article class's default vertical air.
% ---------------------------------------------------------------------------
\titleformat{\section}{\normalfont\large\bfseries}{}{0pt}{}[\vspace{-0.55em}\rule{\linewidth}{0.4pt}]
\titlespacing*{\section}{0pt}{0.85em}{0.35em}
\titleformat{\subsection}{\normalfont\normalsize\bfseries}{}{0pt}{}
\titlespacing*{\subsection}{0pt}{0.6em}{0.2em}
\titleformat{\subsubsection}{\normalfont\small\bfseries\scshape}{}{0pt}{}
\titlespacing*{\subsubsection}{0pt}{0.45em}{0.12em}

% A short run-in heading for card interiors that must not consume a full line.
\newcommand{\telosrubric}[1]{\textbf{#1}\enspace}

% Cross-reference helper. Every Telos project is a set of loose printed sheets,
% so a reference names the sheet rather than a page number.
\newcommand{\sheetref}[1]{\textit{see the \textbf{#1} sheet}}

% Which decision, ADR or sheet governs the paragraph you are reading. A reader
% who disagrees with an instruction should be able to find the reasoning behind
% it rather than having to take it on trust.
\newcommand{\governedby}[1]{{\footnotesize\textsc{governed by #1}}}

% ---------------------------------------------------------------------------
% Quantities. siunitx is not in this TeX Live split, and the guides need only a
% fixed thin space and upright units, so define that directly rather than take
% a package dependency the documented install would not satisfy.
%   \qty{9}{kV}  ->  9 kV        \unitof{kV}  ->  kV
% ---------------------------------------------------------------------------
\newcommand{\unitof}[1]{\ensuremath{\mathrm{#1}}}
\newcommand{\qty}[2]{#1\,\unitof{#2}}
\newcommand{\numrange}[2]{#1\,--\,#2}
\newcommand{\qtyrange}[3]{#1\,--\,#2\,\unitof{#3}}
% Inches and feet appear constantly in the fishing guide; keep one spelling.
\newcommand{\inch}[1]{#1\,in}
\newcommand{\feet}[1]{#1\,ft}

% ---------------------------------------------------------------------------
% Cards and boxes.
% ---------------------------------------------------------------------------

% Titled card with a hairline frame. The default card for facts a reader looks
% up rather than reads through.
\newtcolorbox{teloscard}[1]{
  enhanced, breakable,
  colback=telospaper, colframe=telosink, colbacktitle=telospaper,
  coltitle=telosink,
  boxrule=0.5pt, sharp corners,
  left=6pt, right=6pt, top=4pt, bottom=4pt,
  toptitle=3pt, bottomtitle=3pt, titlerule=0.4pt,
  before skip=6pt, after skip=6pt,
  title={#1}, fonttitle=\small\bfseries
}

% Untitled left-ruled block: hierarchy without a redundant label.
\newtcolorbox{telosframe}{
  enhanced, breakable,
  colback=telospaper, frame hidden, boxrule=0pt, sharp corners,
  borderline west={1.3pt}{0pt}{telosink},
  left=8pt, right=2pt, top=3pt, bottom=3pt,
  before skip=6pt, after skip=6pt
}

% Filled emphasis panel — the "read this first" block. The 10% fill stays
% legible under a cheap laser toner curve and still reads as distinct.
\newtcolorbox{telospanel}[1]{
  enhanced, breakable,
  colback=telosfill, colframe=telosink, colbacktitle=telosfill,
  coltitle=telosink,
  boxrule=0.5pt, sharp corners,
  left=6pt, right=6pt, top=4pt, bottom=4pt,
  toptitle=3pt, bottomtitle=3pt, titlerule=0.4pt,
  before skip=6pt, after skip=6pt,
  title={#1}, fonttitle=\small\bfseries
}

% Hazard block. Deliberately the heaviest object on any page: a 1.4pt frame
% plus a solid title bar. Nothing else in Telos is allowed to look like this,
% so a reader skimming a lesson cannot mistake it for ordinary emphasis.
\newtcolorbox{teloshazard}[1]{
  enhanced, breakable,
  colback=telospaper, colframe=telosink,
  colbacktitle=telosink, coltitle=telospaper,
  boxrule=1.4pt, sharp corners,
  left=6pt, right=6pt, top=5pt, bottom=5pt,
  toptitle=3pt, bottomtitle=3pt,
  before skip=8pt, after skip=8pt,
  title={\MakeUppercase{#1}}, fonttitle=\small\bfseries
}

% A stop-and-check gate. Used before anything destructive, and before anything
% that would be expensive to discover later. Shared by every project that has
% a procedure someone could get hurt following carelessly.
\newtcolorbox{gate}[1]{
  enhanced, breakable,
  colback=telospaper, colframe=telosink,
  boxrule=1.0pt, sharp corners,
  left=6pt, right=6pt, top=4pt, bottom=4pt,
  toptitle=3pt, bottomtitle=3pt, titlerule=0.4pt,
  before skip=7pt, after skip=7pt,
  colbacktitle=telospaper, coltitle=telosink,
  title={\MakeUppercase{#1}}, fonttitle=\small\bfseries
}

% ---------------------------------------------------------------------------
% Rating dots. A five-step scale that prints identically in black-only: filled
% versus open circles, never a shaded bar.
% ---------------------------------------------------------------------------
\newcommand{\telosdot}{\tikz[baseline=-0.42ex]\fill[telosink] (0,0) circle (0.28ex);}
\newcommand{\telosopendot}{\tikz[baseline=-0.42ex]\draw[telosink,line width=0.28pt] (0,0) circle (0.28ex);}
\makeatletter
\newcounter{telos@rate}
\newcommand{\telosrating}[1]{%
  \begingroup
  \setcounter{telos@rate}{0}%
  \@whilenum\value{telos@rate}<#1\do{\telosdot\kern0.22em\stepcounter{telos@rate}}%
  \@whilenum\value{telos@rate}<5\do{\telosopendot\kern0.22em\stepcounter{telos@rate}}%
  \endgroup
}
\makeatother

% ---------------------------------------------------------------------------
% Tables.
%
% tabularx reads its body by scanning for a literal \end{tabularx}, so it can
% never be wrapped in a \newenvironment. Documents therefore open tabularx and
% longtable directly and take their house style from these column types and the
% booktabs rules. The one exception is the fact strip below, which is built on
% plain tabular precisely so it *can* be an environment.
% ---------------------------------------------------------------------------
\newcolumntype{L}{>{\raggedright\arraybackslash}X}
\newcolumntype{R}{>{\raggedleft\arraybackslash}X}
\newcolumntype{C}{>{\centering\arraybackslash}X}
\newcolumntype{B}{>{\bfseries\raggedright\arraybackslash}X}
\newcolumntype{N}{>{\raggedright\arraybackslash}p{0.16\linewidth}}

% Key/value fact strip. Two flush columns of short lookups; the label column is
% a fixed fraction so a stack of cards aligns across the whole guide.
\newenvironment{telosfacts}[1][0.24]
  {\begingroup\small
   \setlength{\parskip}{0pt}%
   \renewcommand{\arraystretch}{1.25}%
   \begin{tabular}{@{}>{\bfseries\raggedright\arraybackslash}p{\dimexpr#1\linewidth\relax}@{\hspace{0.7em}}>{\raggedright\arraybackslash}p{\dimexpr\linewidth-#1\linewidth-0.7em\relax}@{}}}
  {\end{tabular}\par\endgroup}
\newcommand{\telosfact}[2]{#1 & #2\\}

% ---------------------------------------------------------------------------
% Seasonal / diurnal activity bar.
%
% \telosseasonbar{4,3,2,1,1,2,3,3,4,5,4,3} draws a twelve-cell month strip
% where each value 0-5 is a fill height. Height, not gray value, carries the
% signal, so it is exact under any toner curve; the cell is outlined either way
% so an empty month is still visibly a month.
% ---------------------------------------------------------------------------
\newcommand{\telosseasonlabels}{J,F,M,A,M,J,J,A,S,O,N,D}
\newcommand{\telosseasonbar}[1]{%
  \begingroup
  \begin{tikzpicture}[x=1.25em,y=2.1ex,baseline=0pt]
    \foreach \v [count=\i from 0] in {#1} {
      \draw[telosrule,line width=0.3pt] (\i,0) rectangle (\i+0.86,5);
      \ifnum\v>0
        \fill[telosink] (\i,0) rectangle (\i+0.86,\v);
      \fi
    }
    \foreach \m [count=\i from 0] in \telosseasonlabels {
      \node[font=\tiny,anchor=north,inner sep=1.2pt] at (\i+0.43,-0.1) {\m};
    }
  \end{tikzpicture}%
  \endgroup
}

% Four-cell day-part strip: dawn, midday, dusk, night.
\newcommand{\telosdaylabels}{Dawn,Mid,Dusk,Night}
\newcommand{\telosdaybar}[1]{%
  \begingroup
  \begin{tikzpicture}[x=2.9em,y=2.1ex,baseline=0pt]
    \foreach \v [count=\i from 0] in {#1} {
      \draw[telosrule,line width=0.3pt] (\i,0) rectangle (\i+0.92,5);
      \ifnum\v>0
        \fill[telosink] (\i,0) rectangle (\i+0.92,\v);
      \fi
    }
    \foreach \m [count=\i from 0] in \telosdaylabels {
      \node[font=\tiny,anchor=north,inner sep=1.2pt] at (\i+0.46,-0.1) {\m};
    }
  \end{tikzpicture}%
  \endgroup
}

% A bar needs vertical room for its labels when it sits in a fact strip.
\newcommand{\telosbarrow}[1]{\rule{0pt}{4.2ex}#1\rule[-1.6ex]{0pt}{1.6ex}}

% ---------------------------------------------------------------------------
% Plates. A species or apparatus illustration with a caption rule beneath it.
% \telosplate{width fraction}{file}{caption}
% ---------------------------------------------------------------------------
\newcommand{\telosplate}[3]{%
  \begingroup
  \centering
  \includegraphics[width=#1\linewidth]{#2}\par
  \vspace{0.15em}
  \begin{minipage}{#1\linewidth}
    \rule{\linewidth}{0.4pt}\par
    \vspace{0.1em}
    {\footnotesize #3\par}
  \end{minipage}\par
  \endgroup
  \vspace{0.3em}
}

% TikZ drawing conventions shared by every rig, knot, and circuit diagram.
\tikzset{
  telos line/.style={draw=telosink,line width=0.7pt},
  telos hair/.style={draw=telosink,line width=0.35pt},
  telos heavy/.style={draw=telosink,line width=1.2pt},
  telos node/.style={draw=telosink,fill=telospaper,align=center,inner sep=4pt,font=\small},
  telos emphasis/.style={draw=telosink,line width=1.2pt,fill=telospaper,align=center,
    inner sep=4pt,font=\small\bfseries},
  telos label/.style={font=\scriptsize,inner sep=1.5pt},
  telos arrow/.style={-{Stealth[length=1.8mm]},draw=telosink,line width=0.6pt},
  telos dim/.style={|<->|,draw=telosink,line width=0.35pt},
  telos water/.style={draw=telosrule,line width=0.4pt,decorate,
    decoration={snake,amplitude=0.4mm,segment length=3mm}},
}

% ---------------------------------------------------------------------------
% Lesson furniture
% ---------------------------------------------------------------------------

% \lessonhead{number}{time}{who does what}
\newcommand{\lessonhead}[3]{%
  \par\noindent
  \begin{minipage}{\linewidth}
    \rule{\linewidth}{0.4pt}\par
    \vspace{0.3em}
    {\small\textbf{Lesson #1}\quad$\cdot$\quad #2\quad$\cdot$\quad #3\par}
    \vspace{0.25em}
    \rule{\linewidth}{0.4pt}
  \end{minipage}\par
  \vspace{0.5em}}

% The materials list. Two columns of short lines; a lesson that needs a long
% shopping list is a lesson that has not been thought through.
\newenvironment{materials}
  {\begin{telospanel}{What you need}\begin{multicols}{2}
   \begin{itemize}[leftmargin=1em,itemsep=0.05em,topsep=0pt]}
  {\end{itemize}\end{multicols}\end{telospanel}}

% The four rubrics every lesson carries, in the same order every time:
% what you are about to see, what to do, what it means, and what to ask.
\newcommand{\lessonsection}[1]{\section*{#1}\vspace{-0.35em}\rule{\linewidth}{0.4pt}\vspace{0.2em}}

% A prediction box. The point of the curriculum is that they commit to an
% answer before the demonstration, not after.
\newtcolorbox{predict}{
  enhanced, breakable,
  colback=telospaper, colframe=telosink,
  boxrule=0.5pt, sharp corners,
  left=6pt, right=6pt, top=4pt, bottom=4pt,
  before skip=6pt, after skip=6pt,
  borderline west={2.2pt}{0pt}{telosink},
  title={\small\bfseries Predict first}, fonttitle=\small\bfseries,
  colbacktitle=telospaper, coltitle=telosink,
  toptitle=3pt, bottomtitle=3pt, titlerule=0.4pt,
}

% ---------------------------------------------------------------------------
% Colophon. Rendered once, at the end of the last page, in the recovered footer
% area so it never creates a page of its own.
% ---------------------------------------------------------------------------
\newif\ifTelosColophonRendered
\newcommand{\TelosColophonExtra}{}
\newcommand{\TelosColophon}{%
  \ifTelosColophonRendered\else
  \global\TelosColophonRenderedtrue
  \par
  \enlargethispage{0.5in}%
  \vspace{0.3em}%
  \begingroup
  \setlength{\parskip}{0pt}%
  \begin{minipage}{\linewidth}
  \rule{\linewidth}{0.35pt}\par
  \vspace{0.3em}%
  \fontsize{7}{8.2}\selectfont
  \textbf{Telos.} \TelosProject{} — \TelosDocTitle. Revised \TelosRevision.
  \TelosColophonExtra{}
  Project-created text and diagrams are licensed \href{https://creativecommons.org/licenses/by/4.0/}{CC BY 4.0}.
  Incorporated public-domain artwork remains public domain and is credited on
  the plate it appears with. See \texttt{LICENSE} and \texttt{CREDITS.md} in the
  source. Nothing here is professional advice; verify current regulations and
  follow the stated safety procedure.
  \end{minipage}
  \endgroup
  \fi
}
\AtEndDocument{\TelosColophon}

% Shared macros for the homelab project.
%
% This project's documents are read in two very different situations: at a desk
% while designing, and standing in front of a broken machine at eleven at night
% with no network. The second case is the one that sets the style. Commands are
% set apart so they can be found by eye, every dangerous step has a visible
% verification after it, and nothing depends on being able to reach a web page.

% ---------------------------------------------------------------------------
% Decision-record entries, used by the generated digest.
% ---------------------------------------------------------------------------
\newcommand{\adrentry}[4]{%
  \par\vspace{0.7em}
  \needspace{4\baselineskip}%
  \noindent
  \begin{minipage}{\linewidth}
    \rule{\linewidth}{0.6pt}\par
    \vspace{0.25em}
    {\small\textbf{ADR #1}\quad #2\par}
    \vspace{0.1em}
    {\fontsize{7}{8.2}\selectfont\textsc{#3}\quad$\cdot$\quad #4\par}
    \vspace{0.2em}
    \rule{\linewidth}{0.3pt}
  \end{minipage}\par
  \vspace{0.3em}}

% ---------------------------------------------------------------------------
% Procedures.
%
% A rebuild manual is a sequence of commands and the observable evidence that
% each one worked. These environments keep that pairing structural rather than
% a matter of the author remembering.
% ---------------------------------------------------------------------------

% A block of shell to type. Deliberately not a listing package: plain verbatim
% in a ruled box prints identically everywhere and has no font dependency.
\newenvironment{shellblock}
  {\par\vspace{0.25em}\noindent
   \begin{minipage}{\linewidth}
   \rule{\linewidth}{0.3pt}\par\vspace{0.15em}
   \begingroup\footnotesize\ttfamily\obeylines\obeyspaces\parindent=0pt\parskip=0pt}
  {\endgroup\par\vspace{0.15em}\rule{\linewidth}{0.3pt}\end{minipage}\par\vspace{0.25em}}

% What you should see if the previous step worked. A step without one of these
% is a step you cannot verify, and in a rebuild manual that is a defect.
\newcommand{\expectthis}[1]{%
  \par\vspace{0.1em}
  \noindent\begin{minipage}{\linewidth}
  \hspace*{0.6em}\begin{minipage}{\dimexpr\linewidth-0.6em\relax}
  {\footnotesize\textbf{Expect:} #1\par}
  \end{minipage}\end{minipage}\par\vspace{0.3em}}

% \begin{gate} now lives in common/preamble.tex, shared with the other
% projects that document procedures.

% \governedby now lives in common/preamble.tex.

% An instance value that comes from the private overlay. In the published
% documents these render as visible placeholders, which is deliberate: a
% published manual must be complete and obviously generic, never look like it
% is describing somebody's actual network.
\newcommand{\instancevalue}[2]{\texttt{\textlangle #1\textrangle}%
  \footnote{Instance value. Supplied from \texttt{homelab/instance/} at build
  time; #2.}}
\newcommand{\placeholder}[1]{\texttt{\textlangle #1\textrangle}}

% ---------------------------------------------------------------------------
% Role and boundary diagrams.
% ---------------------------------------------------------------------------
\tikzset{
  role/.style={draw=telosink,line width=0.9pt,fill=telospaper,align=center,
    inner sep=5pt,font=\small,minimum width=2.6cm},
  roleoff/.style={draw=telosmute,line width=0.5pt,dash pattern=on 2.5pt off 1.8pt,
    fill=telospaper,align=center,inner sep=5pt,font=\small,minimum width=2.6cm},
  boundary/.style={draw=telosmute,line width=0.5pt,dash pattern=on 3pt off 2pt,
    rounded corners=3pt},
  flow/.style={-{Stealth[length=2mm]},draw=telosink,line width=0.7pt},
  hlabel/.style={font=\fontsize{6.4}{7.4}\selectfont},
}


\telosmeta{Homelab}{Controller Rebuild}
  {Bare metal to serving DHCP, assuming nothing else works}{2026-07-25}

\begin{document}
\telosmasthead[Standalone reconstruction manual]

This manual assumes the Controller is dead and that nothing it hosted is
available: no DNS, no PXE, no artifact service, no Git remote you can reach
through it. It assumes you have this sheet on paper, an existing workstation,
and the recorded inputs for the machine you are rebuilding. If any step here
requires something the Controller was providing, that step is a defect --- report
it rather than working around it.

\begin{gate}{Before you start, confirm you have all five}
\textbf{1.} This document, on paper or on a device that is not the Controller.
\textbf{2.} A working Arch workstation to act as the bootstrap host.
\textbf{3.} The installer image and its published SHA-256, or the ability to
build it from this repository.
\textbf{4.} The recorded instance values for this Controller: hostname, managed
subnet, Controller address, DHCP pool, and the managed interface's MAC address.
\textbf{5.} The LUKS2 recovery material, held off-host. Without it an existing
encrypted disk is unrecoverable, and that is by design.
\end{gate}

\section{Stage 0 --- Prepare the bootstrap host}
\governedby{ADR 0052}

The bootstrap host runs the same dnsmasq and nginx configuration the Controller
profile uses, generated by the same code, so the path you are about to exercise
is the path you have already tested.

\begin{shellblock}
git clone <this repository> telos \&\& cd telos/homelab
python3 -m unittest discover -s tests -t .
\end{shellblock}
\expectthis{\texttt{OK}. If the test suite fails, stop. The installer's
validation logic is what stands between an operator error and a wiped disk, and
it is not optional.}

Build or fetch the installer image and verify it:

\begin{shellblock}
sha256sum -c artifacts/manifest.sha256
\end{shellblock}
\expectthis{Every line reports \texttt{OK}. A mismatch means the artifact is not
the one that was tested --- do not proceed. \governedby{ADR 0043, ADR 0048}}

\begin{gate}{If the network already has a DHCP server}
Run the bootstrap dnsmasq in \textbf{ProxyDHCP} mode so it offers only boot
options and never addresses. Two address-assigning DHCP authorities on one
layer-2 network is forbidden at every stage, including this one.
\governedby{ADR 0008}
\end{gate}

\section{Stage 1 --- Network boot the target}
Set the target's firmware to UEFI network boot. The installer refuses a target
booted in legacy BIOS or CSM mode and stops before offering any destructive
action. \governedby{ADR 0019}

\expectthis{The Archiso environment reaches a shell and the installer starts.
It has made no change to any disk at this point and will not until Stage 3.}

\section{Stage 2 --- Answer, and read the summary}
\governedby{ADR 0004, ADR 0045, ADR 0050}

The installer collects, in order: the profile, the hostname, the target disk,
the managed interface, and --- when Controller network services are enabled --- the
four network inputs. It validates each as you enter it.

\begin{tabularx}{\linewidth}{@{}p{0.26\linewidth}L@{}}
\toprule
\textbf{Input} & \textbf{What is checked}\\
\midrule
Target disk & Confirmed by model, serial and size. Never by \texttt{/dev/sdX}, because that is not stable across boots. Disks under 64 GiB are refused. \governedby{ADR 0051}\\
Managed interface & Chosen from an enumerated list showing kernel name, permanent MAC, link state and speed. There is no default. The MAC becomes the identity and is pinned to \texttt{lan0}.\\
\texttt{managed\_ipv4\_cidr} & Parses as IPv4, and names the network address itself --- a CIDR with host bits set is rejected rather than silently normalised.\\
\texttt{controller\_ipv4\_address} & Usable unicast address inside the subnet, and \emph{outside} the DHCP pool.\\
\texttt{dhcp\_pool\_start/end} & Both usable and inside the subnet; start not greater than end.\\
\bottomrule
\end{tabularx}

Everything else is derived and never prompted for: netmask, network and
broadcast addresses, the client DNS server (the Controller), the suffix
(\texttt{home.arpa}), and the absence of a default-router option.

\begin{gate}{The preflight summary is the authorization boundary}
The summary states the exact disk that will be destroyed, by model and serial;
the profile and hostname; and the complete entered and derived network plan.
Read the disk serial aloud against the label on the drive. Physical access plus
this confirmation \emph{is} the authorization model --- there is no second
approval anywhere. Selecting a profile has authorized nothing.
\end{gate}

\section{Stage 3 --- Install}
After confirmation the installer partitions, encrypts, creates filesystems,
installs the base system with both kernels, writes systemd-boot and the UKIs,
configures \texttt{lan0} and the static address, installs \texttt{sshd} with your
public key, writes the manifest, syncs, and \textbf{powers off}. It does not
reboot. \governedby{ADR 0010}

\expectthis{A clean poweroff, not a reboot. If the machine reboots into the
installed system while still on the provisioning network, power it off before
it can start DHCP.}

\section{Stage 4 --- Relocate, then first boot}
Move the powered-off Controller to the isolated switch and power it on. Being
powered off during the move is the primary DHCP-conflict safety boundary; it is
physical, not a configuration flag.

Under \texttt{development-proof} mode you will be asked for the LUKS2
passphrase at every boot. That is expected --- Milestone A enrolls no TPM token.
\governedby{ADR 0043}

First-boot activation verifies, before starting anything: that \texttt{lan0}
exists, carries the expected MAC and has carrier; that the static address is
configured; and that no other DHCP authority answers on the segment.

\expectthis{\texttt{systemctl is-active dnsmasq nginx} reports \texttt{active}
for both. If activation failed, both are stopped and \texttt{journalctl -u
homelab-first-boot} states which check failed. A partial start is a failure.
\governedby{ADR 0009, ADR 0012}}

\section{Stage 5 --- Prove it}
From a separate client on the isolated switch:

\begin{shellblock}
ip addr show          \# a lease inside the configured pool
ip route show         \# no default route
dig +short \placeholder{hostname}.home.arpa @\placeholder{controller-address}
\end{shellblock}

\begin{telosfacts}[0.30]
\telosfact{Lease}{Inside the pool, and never the Controller's own address.}
\telosfact{No default route}{The client must have none. A default route here means dnsmasq advertised a router that does not exist, and every client on this segment is black-holed. \governedby{ADR 0011}}
\telosfact{DNS}{Returns the Controller's address.}
\telosfact{PXE}{A second target network boots from this Controller.}
\telosfact{Idempotence}{A second Ansible converge run reports no changes.}
\end{telosfacts}

\section{Stage 6 --- Converge}
\governedby{ADR 0053}

Everything beyond the base install is Ansible over SSH from this repository.
Run it in check mode first and read the diff.

\begin{shellblock}
ansible-playbook -i \placeholder{inventory} controller.yml --check --diff
ansible-playbook -i \placeholder{inventory} controller.yml
ansible-playbook -i \placeholder{inventory} controller.yml --check
\end{shellblock}
\expectthis{The third run reports zero changed tasks. A role that is not
idempotent is a role that will surprise you during an outage.}

\section{Recovering instead of rebuilding}
\begin{telosfacts}[0.24]
\telosfact{Bad update}{Boot the recovery UKI, restore the operating-system checkpoint and its matching ESP artifact set, reboot. The checkpoint is taken before the transaction, not after. \governedby{ADR 0028}}
\telosfact{Boot broken}{The recovery UKI is self-contained and boots without the installed root. It starts non-destructive and requires an explicit action before mounting anything writable. \governedby{ADR 0024}}
\telosfact{TPM unlock fails}{Milestone B only. Early boot falls back to the recovery key, which is held off-host. Test that path \emph{before} enabling TPM unlock, never after. \governedby{ADR 0029}}
\telosfact{Disk dead}{Rebuild from this manual. Application data restores from its own backup, on its own schedule --- it is not in the operating-system checkpoint. \governedby{ADR 0027, ADR 0054}}
\telosfact{Recovery key lost}{The data is gone. This is the intended behaviour of encryption at rest, and it is why custody is a decision with its own ADRs rather than an afterthought.}
\end{telosfacts}

\begin{telospanel}{Test this manual by using it}
The QEMU acceptance matrix runs this entire sequence unattended against virtual
hardware and asserts every expectation above. Any change to the installer, the
profile or these instructions must pass it before touching real metal.
A reconstruction manual nobody has followed end to end is a hypothesis.
\governedby{ADR 0056}
\end{telospanel}
\end{document}
